\section*{अध्याय \ref{ch:b3:structural-biology} --- प्रोटीनों का संरचनात्मक जीवविज्ञान}

\begin{solution}{exo:b3:structural-biology:1}
प्राथमिक: $141$ अवशेषों का अनुक्रम, किसी $\alpha$ शृंखला का।
द्वितीयक: उसकी आठ $\alpha$-कुंडलियाँ। तृतीयक: अपने हीम के चारों ओर
लिपटी एक शृंखला का ग्लोबिन \omterm{def:b3:structural-biology:levels}{वलन}। चतुष्क: $\alpha_{2}\beta_{2}$
चतुष्लक और उसके द्विलकों का एक-दूसरे के सापेक्ष घूमने का ढंग।
\end{solution}

\begin{solution}{exo:b3:structural-biology:2}
$36\times \qty{0.15}{nm} = \qty{5.4}{nm}$; $36/3.6 = 10$ फेरे;
$36 - 4 = 32$ हाइड्रोजन बंध।
\end{solution}

\begin{solution}{exo:b3:structural-biology:3}
यह कि प्राकृत संरचना, $105$ संभावनाओं में से चार डाइसल्फ़ाइडों के सही
युग्मन सहित, अवलित शृंखला से स्वतःस्फूर्त रूप से लौट आती है: अनुक्रम
में वह सारी सूचना है जो चाहिए, और प्राकृत अवस्था ऊष्मागतिकीय न्यूनतम
है। शुद्ध प्रोटीन के साथ परखनली में ऐसा करने से कोई भी साँचा, एंज़ाइम
या कोशिकीय तंत्र सूचना के स्रोत के रूप में बाहर हो गया।
\end{solution}

\begin{solution}{exo:b3:structural-biology:4}
NMR: छोटे प्रोटीन, लगभग \qty{40}{kDa} से नीचे, विलयन में।
क्रिस्टलिकी: कोई भी आकार जो क्रिस्टलित हो, पेप्टाइडों से लेकर
राइबोसोमों तक। \omterm{def:b3:structural-biology:methods}{क्रायो-इलेक्ट्रॉन सूक्ष्मदर्शिकी}: बड़े संकुल, लगभग
\qty{50}{kDa} से ऊपर, बिना क्रिस्टलों के। NMR गति सीधे बताता है (और
क्रायो-EM अणुओं को संरूपणी वर्गों में छाँट सकती है); कोई क्रिस्टल
संरचना एक स्थैतिक औसत है।
\end{solution}

\begin{solution}{exo:b3:structural-biology:5}
$2^{150} \approx 1.4\times 10^{45}$ संरूपताएँ, $1.4\times 10^{33}$
s $\approx 5\times 10^{25}$ वर्ष। $3^{20} \approx 3.5\times 10^{9}$,
$3.5$ ms: पेप्टाइड एक सेकंड के भीतर पूरी खोज कर सकता है; प्रोटीन
ब्रह्मांड की आयु में भी नहीं।
\end{solution}

\begin{solution}{exo:b3:structural-biology:6}
$Y = [\alpha(1+\alpha) + Lc\alpha(1+c\alpha)]/[(1+\alpha)^{2} +
L(1+c\alpha)^{2}]$। $\alpha = 1$: $3.01/106 = 0.028$ (एक स्थल
$0.50$)। $\alpha = 10$: $121/242 = 0.50$ ($0.91$)। $\alpha = 100$:
$\num{10300}/\num{10601} = 0.97$ ($0.99$)। MWC वक्र कम लिगंड पर एक
स्थल से बहुत पीछे रहता है, फिर तीव्रता से उठता है --- दो दशकों में
$0.03$ से $0.97$ तक, जहाँ एक स्थल $0.5$ से $0.99$ तक जाता है: वही
तीव्रता \omterm{def:b3:structural-biology:allostery}{सहकारिता} है।
\end{solution}

\begin{solution}{exo:b3:structural-biology:7}
द्विफ़ॉस्फ़ोग्लिसरेट केवल \omterm{def:b3:structural-biology:allostery}{T अवस्था} से बँधता है (केंद्रीय गुहा में),
इसलिए वह T को स्थिर करता है: प्रतिरूप में वह $L$ को, यानी अलिगंडित
प्रोटीन के $T/R$ अनुपात को, ऊपर उठाता है, और इस तरह वक्र को दाईं ओर
खिसका देता है --- कम आसक्ति, ऊतक-दाबों पर अधिक ऑक्सीजन मुक्त। संचित
लाल कोशिकाएँ यह यौगिक खो देती हैं, $L$ गिरता है, वक्र बाईं ओर खिसकता
है, और आधानित रक्त ऑक्सीजन ऊतकों में देने के बजाय पकड़े रहता है, जब तक
कोशिकाएँ एक दिन में उसे फिर से संश्लेषित न कर लें।
\end{solution}

\begin{solution}{exo:b3:structural-biology:8}
$d = \lambda/(2\sin\theta) = 0.1/(2\times 0.259) = \qty{0.19}{nm}$:
आधार-ढाँचा और हर पार्श्व-शृंखला रखी जाती है, और अलग-अलग परमाणु विभेदित
होते हैं। \qty{0.12}{nm} के लिए: $\sin\theta = 0.1/0.24 = 0.417$, $\theta =
\qty{25}{\degree}$।
\end{solution}

\begin{solution}{exo:b3:structural-biology:9}
अध्रुवी पार्श्व-शृंखलाओं के बीच दबा कोई आवेशित लाइसीन दसियों किलोजूल
प्रति मोल की लागत लेता है (उसका आवेश असॉल्वेटित, ल्यूसीन का जलविरोधी
दबाव खोया हुआ), जो \omterm{def:b3:structural-biology:levels}{वलन} के पूरे $\Delta G$ से अधिक है: प्रोटीन अवलित हो
जाता है या क्रोड फिर से जम जाता है। पृष्ठ पर लाइसीन वैसे ही जलयोजित
रहता है जैसे अवलित अवस्था में रहता और कुछ भी नहीं खोता।
$\qty{-30}{kJ/mol} \approx -12RT$ का अर्थ है $K = e^{12} \approx 1.6\times
10^{5}$: किसी भी क्षण \num{160000} में एक अणु अवलित है, और पूरा
अंतराल सैकड़ों में से दो या तीन हाइड्रोजन बंधों के बराबर है --- एक
उत्परिवर्तन जितना।
\end{solution}

\begin{solution}{exo:b3:structural-biology:10}
$R$ स्पीशीज़ का कुल $[R_{0}](1+\alpha)^{n}$ है और $T$ स्पीशीज़ का
$[R_{0}]L(1+c\alpha)^{n}$, इसलिए $\bar R = (1+\alpha)^{n}/[(1+\alpha)^{n}
+ L(1+c\alpha)^{n}]$। $c \to 0$ के लिए $T$ पद $L$ है, और $\bar R =
1/2$ तब होता है जब $(1+\alpha)^{n} = L$, यानी $\alpha = L^{1/n} - 1$।
हीमोग्लोबिन के लिए $(3\times 10^{5})^{1/4} = 23.4$, इसलिए $\alpha \approx 22$:
\qty{22}{mmHg}, जो \qty{26}{mmHg} के $P_{50}$ के निकट है ---
संरूपणी स्विच और अर्ध-संतृप्ति लगभग एक साथ पड़ते हैं।
\end{solution}

\begin{solution}{exo:b3:structural-biology:11}
उसका विरोध इसलिए हुआ कि हर ज्ञात संक्रामक कारक न्यूक्लिक अम्ल के
ज़रिए प्रतिकृत होता था, और कोई प्रोटीन अपना अनुक्रम नकल नहीं कर सकता।
यदि कोई न्यूक्लिक अम्ल आवश्यक पाया जाए, या शुद्ध किया गया प्रोटीन
संचरण में असमर्थ दिखे, तो उसका खंडन हो जाएगा। समर्थन: संक्राम्यता
न्यूक्लिएज़ों, पराबैंगनी और आयनकारी विकिरण की उन मात्राओं को झेल जाती
है जो किसी भी जीनोम को नष्ट कर दें, पर प्रोटीन विकृतकों और प्रोटिएज़ों
से नष्ट हो जाती है; पात्रे तंतुकों में वलित किया गया पुनर्योगज
PrP पशुओं में रोग संचरित करता है, और संचरण पर वही प्रजाति-गुण दिखाता
है। बीज इसलिए चाहिए कि किसी सामान्य एकलक का अकेले बदलना असंभव जितना
धीमा है --- पहले अनुप्रस्थ-$\beta$ संरचना का नाभिक बनना चाहिए, जो
दर-निर्धारक घटना है --- जबकि कोई पहले से बना तंतुक-सिरा एकलकों को तेज़ी
से बदल देता है: आकार साँचे से बढ़कर नकल होता है।
\end{solution}

\begin{solution}{exo:b3:structural-biology:12}
अनेक अनुक्रम एक ही संरचना में वलित होते हैं: किसी \omterm{def:b3:structural-biology:levels}{वलन} को चाहिए
जलविरोधी और ध्रुवी स्थितियों का कोई प्रतिरूप, मोड़ों पर कुछ ग्लाइसीन
और प्रोलीन, और वे अवशेष जो काम करते हैं; पृष्ठ पर बाकी सब बिना किसी
दंड के बदल सकता है, इसलिए वहाँ उत्परिवर्तन जमा होते जाते हैं जबकि चयन
\omterm{def:b3:structural-biology:levels}{वलन} और प्रकार्य बचाए रखता है। इसलिए संरक्षित अवशेष हैं क्रोड (पहचानों
से अधिक किसी प्रतिरूप के रूप में), उत्प्रेरक तथा बंधक अवशेष, और
संरचनात्मक ग्लाइसीन, प्रोलीन तथा सिस्टीन। प्रोफ़ाइल विधियाँ ठीक उन्हीं
स्तंभों को भार देती हैं और इसलिए \qty{15}{\%} समरूपता पर भी कुल के
सदस्यों को पहचान लेती हैं; प्रोटीन-अभिकल्प इसी तर्क को उलटा चलाता है,
किसी लक्ष्य \omterm{def:b3:structural-biology:levels}{वलन} के लिए कोई जलविरोधी प्रतिरूप चुनकर और पृष्ठ को कुछ भी
रहने देकर।
\end{solution}

\begin{solution}{pb:b3:structural-biology:1}
\textbf{1.} $0.75\times 150 \approx 112$ कुंडलित अवशेष; $112\times
\qty{0.15}{nm} = \qty{17}{nm}$; $112/3.6 \approx 31$ फेरे।
\textbf{2.} $112 - 8\times 4 = 80$ हाइड्रोजन बंध।
\textbf{3.} द्रव्यमान $150\times 110 = \qty{16500}{Da} = 2.7\times 10^{-20}$
g; आयतन $2.7\times 10^{-20}\times 0.73 = 2.0\times 10^{-20}$ cm$^{3}$
$= \qty{20}{nm^3}$; त्रिज्या $(3V/4\pi)^{1/3} = \qty{1.7}{nm}$।
\textbf{4.} फैली हुई $150\times 0.36 = \qty{54}{nm}$, \qty{3.4}{nm}
के व्यास के सामने: $16$ का गुणक।
\textbf{5.} कुंडलित अवशेष फैली हुई लंबाई के \qty{75}{\%} हैं,
$112\times 0.36 = \qty{40}{nm}$; कुंडली उन्हें
\qty{17}{nm} तक दबा देती है, गुणक $0.36/0.15 = 2.4$।
\textbf{6.} $40\times 4 = \qty{160}{kJ/mol}$ का लाभ, $150\times 0.9 =
\qty{135}{kJ/mol}$ का ह्रास: शुद्ध $\Delta G \approx \qty{-25}{kJ/mol}$,
सीमांत परास में।
\textbf{7.} $3^{150} \approx 4\times 10^{71}$; $4\times 10^{59}$ s
$\approx 10^{52}$ वर्ष।
\textbf{8.} $10^{-2}\times 10^{12} = 10^{10}$ संरूपताएँ, गिनती का
$3\times 10^{-62}$ अंश।
\textbf{9.} कुंडलियाँ \qty{1}{\micro s} में (समांतर में); $8! = \num{40320}$
जमाव प्रति सेकंड $10^{6}$ की दर से \qty{40}{ms} लेते हैं: कुल लगभग
\qty{40}{ms}, यानी प्रेक्षित कोटि। हाँ: भागों को स्वतंत्र रूप से और
फिर पूरे को हल करना एक कीप है --- खोज-समष्टि गुणनखंडित है, पूरी छानी
नहीं जाती।
\textbf{10.} प्रत्याशित चक्र $1/0.3 = 3.3$, लगभग \qty{33}{s}; छह
चक्रों के बाद $0.7^{6} = 0.12$ अवलित बचता है।
\textbf{11.} $\Delta G = -RT\ln K$ के साथ $K = 10^{4} \to 10^{2}$:
$-23.7$ से \qty{-11.9}{kJ/mol} तक, यानी \qty{12}{kJ/mol} स्थायित्व का
ह्रास ($RT = \qty{2.58}{kJ/mol}$, $\ln 100 = 4.6$)।
\textbf{12.} पुंजन किसी नाभिक से आरंभ होता है, जिसके बनने की दर
पुंजन-प्रवण अवलित स्पीशीज़ की सांद्रता की किसी ऊँची घात की तरह बढ़ती
है; उस स्पीशीज़ का सौ गुना होना किसी जीवन-काल के भीतर बीज बनने की
संभावना को कई कोटियों तक बढ़ा देता है, और बीज बन जाने पर वृद्धि तेज़
तथा स्वयं-साँचा होती है।
\textbf{13.} $\alpha = 100$: $Y = 1.054\times 10^{8}/1.089\times 10^{8} =
0.97$। $\alpha = 20$: अंश $185\,220 + 103\,680 = 288\,900$,
हर $194\,481 + 622\,080 = 816\,561$: $Y = 0.35$।
\textbf{14.} फेफड़े में $\bar R = 1.041\times 10^{8}/1.089\times 10^{8} = 0.96$;
पेशी में $194\,481/816\,561 = 0.24$।
\textbf{15.} क्षमता का $0.97 - 0.35 = 0.62$ उतरता है (भार का
दो-तिहाई)। एक स्थल: $100/126 - 20/46 = 0.79 - 0.43 = 0.36$।
\textbf{16.} क्षमता $150\times 1.34 = \qty{200}{mL/L}$; पेशी तक
$0.62\times 200 \approx \qty{125}{mL/L}$। विश्राम में $(0.97 - 0.78)\times
200 = \qty{38}{mL/L}$; $250/38 \approx \qty{6.5}{L/min}$ रक्त-प्रवाह
--- विश्राम के हृद्-निर्गम की कोटि।
\textbf{17.} $L = 3\times 10^{6}$, $\alpha = 40$: अंश $2.76\times
10^{6} + 3.29\times 10^{6} = 6.05\times 10^{6}$, हर $2.83\times
10^{6} + 11.5\times 10^{6} = 14.4\times 10^{6}$: $Y = 0.42$, $0.78$ के
बजाय --- विश्राम में कहीं अधिक ऑक्सीजन उतरती है, यही ऊँचाई के लिए रक्त
का और काम के लिए पेशी का अनुकूलन है।
\textbf{18.} भ्रूणीय, $L = 3\times 10^{4}$, $\alpha = 30$: अंश
$893\,730 + 19\,770 = 913\,500$, हर $923\,520 + 85\,680 =
1\,009\,200$: $Y = 0.91$। उसी दाब पर माता का: अंश
$893\,730 + 197\,730 = 1\,091\,460$, हर $923\,520 + 856\,830
= 1\,780\,350$: $Y = 0.61$। समान दाब पर भ्रूणीय रक्त अधिक
ऑक्सीजन रखता है, इसलिए ऑक्सीजन अपरा के पार माता से भ्रूण की ओर बहती है।
\textbf{19.} $(10^{5}\,\text{nm}/6\,\text{nm})^{3} = 4.6\times 10^{12}$
कोष्ठिकाएँ; $1.9\times 10^{13}$ अणु।
\textbf{20.} $d = 0.1/(2\sin 14.5^{\circ}) = \qty{0.20}{nm}$: हाँ,
हर पार्श्व-शृंखला और परमाणु रखा जाता है।
\textbf{21.} $\sin\theta = 0.1/0.3$, $\theta = \qty{19.5}{\degree}$।
किसी अव्यवस्थित क्रिस्टल में अणु ठीक-ठीक पंक्ति में नहीं होते, ऊँचे
कोण पर प्रकीर्ण तरंगें --- जो सूक्ष्म विवरण कूटती हैं --- अब कला में
नहीं जुड़तीं, और धब्बे ऊँचे कोण पर मिट जाते हैं; केवल \omterm{def:b3:structural-biology:methods}{कम-विभेदन} वाले
धब्बे बचते हैं।
\textbf{22.} $d$ को आधा करने के लिए दोगुना संकेत चाहिए, यानी चार गुना
कण: लगभग \num{40000} (व्यवहार में अधिक, क्योंकि दूसरी त्रुटियाँ आ जाती हैं)।
\textbf{23.} झिल्ली प्रोटीन, जो अपमार्जक विलयनों से शायद ही क्रिस्टलित
होते हैं और जिन्हें किसी लिपिड नैनोचक्रिका में चित्रित किया जा सकता
है; बड़े लचीले संकुल --- राइबोसोम, विच्छेदन-काय, अपने नियामकों सहित
आयन वाहिकाएँ --- जिनकी विषमता क्रिस्टलन रोकती है पर जिनके कण संरूपणी
वर्गों में छाँटे जा सकते हैं।
\textbf{24.} बँधा हीम और ऑक्सीजन तथा उनकी ठीक ज्यामिति, क्रमित जल और
आयन, दूसरी संरूपता (T बनाम R) और यह प्रमाण कि \omterm{def:b3:structural-biology:levels}{वलन} सही है --- कोई
पूर्वानुमान बिना किसी प्रायोगिक जाँच के एक अकेला अलिगंडित प्रतिरूप है।
\textbf{25.} क्षमता का $0.62$ उतरना (एक स्थल $0.36$); लेविंथल-काल
लगभग $10^{52}$ वर्ष; \qty{0.20}{nm} पर कोई मानचित्र।
\end{solution}
